\section{Largest Rectangle in Histogram}
\label{sec:sq-largest-rectangle-histogram}

\problemstatement{Given an array $\textit{heights}$ representing the
heights of adjacent bars of width $1$ in a histogram, find the area of
the \textbf{largest rectangle} that can be formed using contiguous bars
(the rectangle's height is limited by the shortest bar it spans).}

\begin{patternbox}
This is the contribution-counting idea from
Section~\ref{sec:sq-sum-subarray-minimums}, repurposed to
\textbf{maximize} rather than sum: for every bar $i$, the largest
rectangle that uses $\textit{heights}[i]$ as its limiting height spans
exactly the maximal contiguous range around $i$ over which
$\textit{heights}[i]$ is the minimum -- found by the exact same nearest
strictly-smaller-element boundaries on each side.
\end{patternbox}

\subsection*{Understanding the problem carefully}

For a rectangle of height $\textit{heights}[i]$ to be valid, every bar it
spans must be at least $\textit{heights}[i]$ tall (otherwise the
rectangle would poke above some shorter bar within its width). So the
widest rectangle using $\textit{heights}[i]$ as its height extends left
until hitting a strictly shorter bar (or the histogram's edge), and right
until hitting a strictly shorter bar (or the edge) -- exactly the
$\textit{left}[i]$ and $\textit{right}[i]$ boundaries from
Section~\ref{sec:sq-sum-subarray-minimums}, just used to compute one
candidate area per bar and take the maximum, rather than summing
contributions.

\begin{ideabox}[One Pop, One Area, via a Single Monotonic Pass]
Maintain a stack of indices with strictly increasing heights. Scan left
to right (appending a sentinel height of $0$ at the very end to force
every remaining bar to be resolved). When the current height is less than
the stack's top height, pop that top index $j$: the rectangle with height
$\textit{heights}[j]$ has width $i - (\text{new top index}) - 1$ (or
simply $i$ if the stack becomes empty, meaning $j$'s rectangle could
extend all the way to the histogram's left edge). Compute this area and
update the running maximum. Continue popping until the current height no
longer triggers a pop, then push the current index.
\end{ideabox}

\subsection*{Why a single left-to-right pass (with a sentinel) suffices}

This combines the boundary-finding idea from
Section~\ref{sec:sq-sum-subarray-minimums} with the floor-resolution idea
from Trapping Rain Water (Section~\ref{sec:sq-trapping-rain-water}): each
bar's rectangle is fully resolved -- both its left and right boundary
determined -- in a single pop, because by the time bar $j$ is popped, the
current index $i$ is exactly its right boundary (the nearest strictly
shorter bar to the right), and the stack's new top after the pop is
exactly its left boundary (the nearest still-standing, and therefore
strictly shorter or absent, bar to the left). The sentinel height $0$ at
the end guarantees every bar still on the stack when the scan finishes is
popped and resolved, rather than left unresolved.

\subsection*{Algorithm}

\begin{enumerate}
  \item Append a sentinel value $0$ to the end of $\textit{heights}$
        (conceptually); initialize an empty stack and
        $\textit{maxArea} \gets 0$.
  \item For $i$ from $0$ to $n$ (inclusive of the sentinel position):
  \begin{enumerate}
    \item While the stack is not empty and
          $\textit{heights}[i] < \textit{heights}[\text{stack's top}]$:
    \begin{itemize}
      \item Pop the top as $j$.
      \item $\textit{width} \gets$ (stack empty $? i :$
            $i - \text{new top} - 1$).
      \item $\textit{maxArea} \gets \max(\textit{maxArea},
            \textit{heights}[j] \times \textit{width})$.
    \end{itemize}
    \item Push $i$.
  \end{enumerate}
  \item Return \textit{maxArea}.
\end{enumerate}

\subsection*{Dry run}

\dryrun{Take $\textit{heights} = [2,1,5,6,2,3]$ (sentinel $0$ appended at
index $6$).}

\begin{itemize}
  \item $i=0$ ($h=2$): stack empty, push. Stack $=[0]$.
  \item $i=1$ ($h=1$): $1<2$, pop $0$. Stack empty: width $=1$. Area
        $=2\times1=2$. $\textit{maxArea}=2$. Push $1$. Stack $=[1]$.
  \item $i=2$ ($h=5$): $5 \not< 1$, no pop. Push. Stack $=[1,2]$.
  \item $i=3$ ($h=6$): $6 \not< 5$, no pop. Push. Stack $=[1,2,3]$.
  \item $i=4$ ($h=2$): $2<6$, pop $3$: width $=4-2-1=1$. Area
        $=6\times1=6$. $\textit{maxArea}=6$. $2<5$, pop $2$: width
        $=4-1-1=2$. Area $=5\times2=10$. $\textit{maxArea}=10$. $2
        \not< 1$, stop. Push $4$. Stack $=[1,4]$.
  \item $i=5$ ($h=3$): $3 \not< 2$, no pop. Push. Stack $=[1,4,5]$.
  \item $i=6$ ($h=0$, sentinel): $0<3$, pop $5$: width $=6-4-1=1$. Area
        $=3$. $0<2$, pop $4$: width $=6-1-1=4$. Area $=8$. $0<1$, pop
        $1$: stack empty: width $=6$. Area $=6$. None exceed $10$.
\end{itemize}

The largest rectangle has area $\boxed{10}$ (bars of height $5$ and $6$
at indices $2,3$, using height $5$ across width $2$).

\subsection*{C++ implementation}

\begin{lstlisting}
#include <bits/stdc++.h>
using namespace std;

int largestRectangleArea(vector<int>& heights) {
    heights.push_back(0); // sentinel
    int n = heights.size();
    stack<int> st;
    long long maxArea = 0;

    for (int i = 0; i < n; i++) {
        while (!st.empty() && heights[i] < heights[st.top()]) {
            int j = st.top();
            st.pop();
            int width = st.empty() ? i : i - st.top() - 1;
            maxArea = max(maxArea, (long long)heights[j] * width);
        }
        st.push(i);
    }
    heights.pop_back(); // restore original array
    return maxArea;
}

int main() {
    vector<int> heights = {2, 1, 5, 6, 2, 3};
    cout << "Largest rectangle area: "
         << largestRectangleArea(heights) << endl;
    return 0;
}
\end{lstlisting}

\begin{complexitybox}
\textbf{Time Complexity.} Each index is pushed once and popped at most
once, giving
\[
  O(n).
\]
\textbf{Space Complexity.} $O(n)$ for the stack.
\end{complexitybox}

\begin{takeawaybox}
``Largest rectangle/area limited by the shortest bar within some span''
is solved by the same boundary-finding idea as
Section~\ref{sec:sq-sum-subarray-minimums}, condensed into a single
left-to-right pass with a sentinel, since we only need each bar's
\emph{maximum} possible area, not a sum over every possible sub-span.
This single-pass technique is reused directly in Maximal Rectangle
(Section~\ref{sec:sq-maximal-rectangle}) to extend from a 1D histogram to
a full 2D binary matrix.
\end{takeawaybox}
